\chapter{Améliorations et suite du projet}
\label{ch:ameliorations}

Le chapitre précédent a listé quelques points encore incertains ou non testés
(voir la synthèse en fin de chapitre~8). Mais l'essentiel est ailleurs~:
passer d'un prototype à un seul poste à une machine de production à 10
postes demande surtout des évolutions structurelles qui n'avaient pas besoin
d'être réglées pour valider le concept, et qui deviennent maintenant
incontournables. C'est l'objet de ce chapitre.

\section{Passage du prototype à une machine à 10 postes}

\subsection{Structure porteuse et socle commun}

Le prototype repose sur un socle individuel, avec ses propres poteaux pour
la caméra et l'éclairage. Le reproduire tel quel dix fois, comme envisagé au
chapitre précédent pour l'estimation d'encombrement, multiplierait d'autant
le nombre de poteaux à fabriquer et aligner, et laisserait peu de place
autour de chaque mouvement pour le fixer.

Un socle principal unique, supportant les 10 modules, serait préférable~: en
mutualisant une partie de la structure porteuse (des traverses communes
plutôt qu'un jeu complet par poste), on réduit le nombre d'éléments verticaux
et on libère de l'espace autour de chaque posage pour la fixation du
mouvement.

\subsection{Gestion des câbles et des alimentations}

Sur le prototype, les câbles moteur, caméra et éclairage cheminent librement
autour de la structure, sans organisation particulière. À l'échelle de 10
postes, ça deviendrait vite un problème, pour la maintenance comme pour la
sécurité~: câbles qui pendent, risques d'accrochage. Le socle principal
devra donc prévoir des canaux pour faire passer proprement ces câbles
jusqu'aux alimentations, ainsi que des emplacements pour fixer les blocs
d'alimentation eux-mêmes, aujourd'hui simplement posés à côté du prototype.

\subsection{Connectivité~: bus USB}

Côté moteurs, pas de souci~: le bus RS485 (tableau~\ref{tab:specs_driver})
permet déjà de connecter jusqu'à 255 modules sur un seul câble, via un
unique adaptateur USB. Les caméras, en revanche, ont chacune leur propre
liaison~: relier directement 10 caméras, et le reste des périphériques par
poste, au même ordinateur reviendrait à brancher une vingtaine de ports USB,
ce qui n'est pas gérable en l'état. Il faudra donc mettre en place un vrai
système de bus USB, par exemple via des concentrateurs industriels, pour
regrouper ces connexions.

\subsection{Adaptations logicielles}

Cette connectivité ne suffira pas à elle seule~: le logiciel actuel ne
pilote pas plusieurs modules en parallèle. Une fois le bus USB en place, le
code devra donc être adapté pour reconnaître et piloter chaque poste
indépendamment, en s'appuyant sur la structure de cases déjà prévue dans
l'interface pour représenter plusieurs modules.

\section{Amélioration du système de fixation du mouvement}

Une fois la machine terminée, il a été remarqué lors de tests avec un mouvement équipé de sa masse oscillante, que celle-ci
rendait l'épaisseur du mouvement inconstante et par conséquent plus compliqué à fixer.
Une solution a alors été imaginée pour pallier ce problème. L'idée est de fixer le mouvement
directement dans sa boîte en plexiglas, plutôt que de le serrer seul,
ce qui demanderait de revoir la cotation de certaines pièces afin de s'assurer que toute l'architecture du système reste compatible et que ça ne désalignerait pas la tige par rapport à l'accouplement.

De plus, cette solution garantirait qu'aucun serrage ne soit exercé sur le mouvement
lui-même.

Toutefois, des cales sous la boîte restent malgré tout nécessaires. En effet, l'analyse de la
saccade a montré qu'une variation de hauteur d'un mouvement à l'autre
désaligne l'accouplement, ce qui induit une erreur à la fois sur la
vitesse et sur la précision de la mesure. Le but de ces cales est que la tige de tous les mouvements soit parfaitement alignée avec l'arbre moteur. Il y aurait donc un modèle de cale à concevoir pour chaque type de mouvement.

\section{Synthèse}

Le tableau~\ref{tab:synthese_ameliorations} résume les principaux points
restant à traiter pour passer du prototype actuel à une machine conforme à
l'ensemble du cahier des charges.

\begin{table}[H]
  \centering
  \caption{Synthèse des points restant à traiter}
  \label{tab:synthese_ameliorations}
  \begin{tabular}{|p{4cm}|p{10.5cm}|}
    \hline
    \textbf{Point à traiter} & \textbf{Description} \\
    \hline
    Structure à 10 postes & Concevoir un socle principal commun, réduisant le
    nombre de poteaux et libérant de l'espace autour de chaque posage. \\
    \hline
    Câblage et alimentations & Ajouter des canaux de câbles et des
    emplacements de fixation dédiés aux alimentations. \\
    \hline
    Connectivité USB & Mettre en place un bus USB pour gérer les caméras des
    10 postes sans multiplier les ports. \\
    \hline
    Logiciel multi-postes & Adapter le code pour piloter chaque poste
    indépendamment via le bus USB. \\
    \hline
    Fixation du mouvement & Fixer le mouvement dans sa boîte en plexiglas
    plutôt que de le serrer seul, en adaptant la cotation des pièces
    concernées et en prévoyant des cales pour préserver l'alignement de
    l'accouplement. \\
    \hline
  \end{tabular}
\end{table}

Deux points supplémentaires n'ont volontairement pas été inclus dans ce
tableau~: contrairement aux points précédents, ils ne sont pas indispensables
pour que la machine réponde au cahier des charges, mais restent conseillés
pour gagner en confiance dans les résultats. Réaliser un plus grand nombre
d'essais en conditions réelles permettrait de réellement garantir la
fiabilité de la mesure, plutôt que de s'appuyer sur les quelques tests
effectués jusqu'ici. De même, une mesure directe du couple, via un banc de
test dédié, offrirait une vérification plus rigoureuse que les calculs
théoriques et les essais qualitatifs réalisés à ce jour. Ni l'un ni l'autre
n'est indispensable, mais avoir de vrais essais pour appuyer ces valeurs
resterait un gage de qualité.
